\subsection{Hessenberg Reduction and QR Eigenvalue Methods}

Throughout this subsection, $n\in\mathbb{Z}_{>0}$ and $\mathbb{F}\in\{\mathbb{R},\mathbb{C}\}$.

\begin{definition} \label{def:upper-hessenberg-matrix} \hlt{Upper Hessenberg Matrix}\\
Matrix $A=[a_{ij}]\in\operatorname{Mat}(n,n,\mathbb{F})$ is \hlt{upper Hessenberg} iff $(\forall\,(i,j)\in J_n^2)\,[i>j+1\Longrightarrow a_{ij}=0]$.
\end{definition}

Numerical eigenvalue methods introduce structural zeros by unitary similarity.\\
Upper Hessenberg form is the standard preliminary reduction for QR iteration.\\
Figure \ref{fig:upper-hessenberg-pattern} separates unrestricted entries from forced zeros below the first subdiagonal.

\begin{illustrationfigure}[H]
  \begin{tikzpicture}[font=\small]
  \node at (0,0) {$
  \begin{bmatrix}
  * & * & * & * & *\\
  * & * & * & * & *\\
  0 & * & * & * & *\\
  0 & 0 & * & * & *\\
  0 & 0 & 0 & * & *
  \end{bmatrix}
  $};
  \draw[red!70, thick] (-1.35,-0.45) -- (0.65,-1.15);
  \node[red!70!black, align=center] at (2.45,-1.0) {first\\subdiagonal};
  \draw[->, red!70, thick] (1.65,-0.95) -- (0.72,-0.92);

  \node[blue!70!black, align=center] at (0,1.75) {entries below this subdiagonal are zero};
  \draw[->, blue!70, thick] (-0.2,1.45) -- (-1.0,-0.85);
\end{tikzpicture}

  \caption{Upper Hessenberg pattern: asterisks are unrestricted and entries below the first subdiagonal are zero.}
  \label{fig:upper-hessenberg-pattern}
\end{illustrationfigure}

\begin{theorem} \label{thm:unitary-upper-hessenberg-reduction} \hlt{Unitary Upper Hessenberg Reduction}\\
Every $A\in\operatorname{Mat}(n,n,\mathbb{F})$ is unitarily similar to an upper Hessenberg matrix with real nonnegative subdiagonal entries.
\begin{proof}
For $n=1$, take $U=I_1$. Assume $n\geq2$ and set $A_0\defeq A$.\\
For $k\in J_{n-1}$ in increasing order, assume columns $1,\ldots,k-1$ have the required form.\\
This condition is vacuous for $k=1$. Let $\mathbf{w}_k\in\mathbb{F}^{n-k}$ be the tail of column $k$ of $A_{k-1}$.\\
If $\mathbf{w}_k=\mathbf{0}$, set $V_k\defeq I_{n-k}$. Otherwise set $\mathbf{v}_k\defeq\frac{\mathbf{w}_k}{\norm{\mathbf{w}_k}_2}$; then $\norm{\mathbf{v}_k}_2=1$.\\
Choose unitary $V_k$ with $V_k\mathbf{e}_1=\mathbf{v}_k$ by Corollary \ref{cor:householder-unitary-annihilator}(iii) in the nonzero case.\\
In either case, $V_k^*\mathbf{w}_k=\norm{\mathbf{w}_k}_2\mathbf{e}_1$ by the defining case and unitarity. Set $U_k\defeq I_k\oplus V_k$ and $A_k\defeq U_k^*A_{k-1}U_k$.\\
The identity $U_k^*U_k=I_n$ holds by block multiplication. The first $k-1$ columns are fixed by right multiplication with $U_k$.\\
They are fixed by left multiplication with $U_k^*$ because only their zero tails are mixed.\\
Column $k$ has the required zeros and subdiagonal entry $\norm{\mathbf{w}_k}_2\in\mathbb{R}_{\geq0}$ by the tail identity.\\
For $U\defeq U_1\cdots U_{n-1}$, $U$ is unitary by identities $U_k^*U_k=I_n$. Identity $A_{n-1}=U^*AU$ holds by iteration.\\
Matrix $A_{n-1}$ has the required form by induction.\\
\hlt{Structural Commentary}: Trailing block similarities preserve completed columns while aligning the next column tail.\\
\hlt{Key Insight}: $V_k^*\mathbf{w}_k=\norm{\mathbf{w}_k}_2\mathbf{e}_1$ installs both zeros and phase normalization.
\end{proof}
\end{theorem}

\begin{perspective} \label{per:hessenberg-before-qr} \hlt{Why Hessenberg Form Matters}\\
Unitary similarity preserves $p_A$ by Theorem \ref{thm:similarity-preserves-characteristic-polynomial}(i), and hence the spectrum by Theorem \ref{thm:matrix-characteristic-polynomial-basic-coefficients}(ii).\\
Upper Hessenberg form leaves only the first subdiagonal to eliminate in a QR factorisation.
\end{perspective}

\begin{computationmethod} \label{meth:qr-eigenvalue-algorithm-preview} \hlt{QR Eigenvalue Algorithm Preview}\\
Let $A\in\operatorname{Mat}(n,n,\mathbb{F})$. Choose unitary $U$ such that $A_0\defeq U^*AU$ is upper Hessenberg by Theorem \ref{thm:unitary-upper-hessenberg-reduction}. For $k\in\mathbb{Z}_{\geq0}$, perform:
\begin{enumerate}[label=(\roman*)]
\setlength{\itemsep}{0pt}
\item Choose $\mu_k\in\mathbb{F}$ and factor $A_k-\mu_kI_n=Q_kR_k$ by Theorem \ref{thm:qr-factorization}(i). Here $Q_k$ is unitary and $R_k$ is upper triangular.
\item Set $A_{k+1}=R_kQ_k+\mu_kI_n=Q_k^*A_kQ_k$.
\end{enumerate}
$(\forall\,k\in\mathbb{Z}_{\geq0})\,A_k\sim A$ by induction through the initialization and update similarities.\\
$(\forall\,k\in\mathbb{Z}_{\geq0})\,\operatorname{spec}(A_k)=\operatorname{spec}(A)$ by the preceding relation, Theorem \ref{thm:similarity-preserves-characteristic-polynomial}(i), and Theorem \ref{thm:matrix-characteristic-polynomial-basic-coefficients}(ii).\\
No shift rule, deflation test, stopping criterion, or convergence claim is specified.
\end{computationmethod}

\begin{computationmethod} \label{meth:hessenberg-qr-step-plane-rotations} \hlt{Hessenberg QR Step by Plane Unitaries}\\
Assume $n\geq2$. Let $H\in\operatorname{Mat}(n,n,\mathbb{C})$ be upper Hessenberg, and let $\mu\in\mathbb{C}$.
\begin{enumerate}[label=(\roman*)]
\setlength{\itemsep}{0pt}
\item Set $B_0\defeq H-\mu I_n$.
\item For $j\in J_{n-1}$, set $\mathbf{z}_j\defeq\left(\begin{smallmatrix}(B_{j-1})_{jj}\\(B_{j-1})_{j+1,j}\end{smallmatrix}\right)$. If $\mathbf{z}_j=\mathbf{0}$, set $\widehat{G}_j\defeq I_2$. Otherwise set $\mathbf{u}_j\defeq\frac{\mathbf{z}_j}{\norm{\mathbf{z}_j}_2}$, so $\norm{\mathbf{u}_j}_2=1$.\\
In the nonzero case, choose unitary $\widehat{G}_j$ with $\widehat{G}_j\mathbf{e}_1=\mathbf{u}_j$ by Corollary \ref{cor:householder-unitary-annihilator}(iii).\\
In either case, $\widehat{G}_j^*\mathbf{z}_j=\norm{\mathbf{z}_j}_2\mathbf{e}_1$ by the construction.\\
Let $G_j$ act as $\widehat{G}_j$ on coordinates $j,j+1$ and as the identity elsewhere. Set $B_j\defeq G_j^*B_{j-1}$.
\item Set $Q\defeq G_1\cdots G_{n-1}$ and $R\defeq B_{n-1}$.
\item Define $H_+\defeq RQ+\mu I_n=Q^*HQ$.
\end{enumerate}
Matrix $B_0$ is upper Hessenberg. At iteration $j$, the two mixed rows vanish in columns $1,\ldots,j-1$ by induction.\\
Their mixing preserves these zeros, and the tail identity annihilates entry $(j+1,j)$.\\
Thus $B_j$ is upper Hessenberg with its first $j$ columns upper triangular by induction.\\
Hence $R$ is upper triangular and $H-\mu I_n=QR$ by $B_{n-1}=Q^*(H-\mu I_n)$.\\
Set $Q_0\defeq I_n$ and $Q_j\defeq G_1\cdots G_j$. Each $Q_{j-1}$ acts only on the first $j$ coordinates.\\
Thus $Q_{j-1}\mathbf{e}_{j+1}=\mathbf{e}_{j+1}$, and $Q_j$ is upper Hessenberg by induction.\\
Matrix $Q$ is unitary by the unitarity of each $G_j$. The identity $H_+=Q^*HQ$ holds by the shifted factorisation.\\
For $i>j+1$, $(RQ)_{ij}=\sum\limits_{\ell=1}^{n}r_{i\ell}q_{\ell j}=0$ by $r_{i\ell}\neq0\Longrightarrow\ell\geq i$ and $q_{\ell j}\neq0\Longrightarrow\ell\leq j+1$. Hence $H_+=RQ+\mu I_n$ is upper Hessenberg.\\
The identity $\operatorname{spec}(H_+)=\operatorname{spec}(H)$ holds by Theorem \ref{thm:similarity-preserves-characteristic-polynomial}(i) and Theorem \ref{thm:matrix-characteristic-polynomial-basic-coefficients}(ii).\\
The construction is exact-arithmetic; no floating-point stability or deflation claim is made.
\end{computationmethod}
